\documentclass[12pt,a4paper]{article}

\usepackage[utf8]{inputenc}
\usepackage[T1]{fontenc}
\usepackage[margin=2.5cm]{geometry}
\usepackage{hyperref}
\usepackage{microtype}
\usepackage{parskip}
\usepackage{lmodern}
\usepackage{amsmath}

\hypersetup{colorlinks=true, urlcolor=blue!60!black}

\pagestyle{empty}

\begin{document}

\begin{flushright}
Kundai Farai Sachikonye \\
Auf der Lichtung 19, 82178 Puchheim, Germany \\
\href{https://github.com/fullscreen-triangle}{github.com/fullscreen-triangle} \\
\texttt{kundai.sachikonye@bitspark.com} \\
(+49) 1626386344 \\[6pt]
18 May 2026
\end{flushright}

\vspace{1em}

Forschungszentrum Jülich GmbH \\
HDS-LEE Graduate School \\
52428 Jülich, Germany \\
Ref.\ 2025D-179

\vspace{1.5em}

\textbf{Re: PhD Position --- Probabilistic Metabolic Flux Analysis (Ref.\ 2025D-179)}

\vspace{1em}

Dear Hiring Committee,

I am applying for the doctoral position in probabilistic metabolic flux analysis
within the HDS-LEE graduate school (Ref.\ 2025D-179).
The position's core challenge --- developing a Bayesian inference framework
robust to model misspecification for hybrid model- and data-driven metabolism
modelling --- sits at the intersection where I have worked most intensively:
hierarchical Bayesian models, differentiable probabilistic inference, and
large-scale biological data, with the specific application domain being
metabolic measurement data from my doctoral research in Computational
Lipidomics at TUM.

\textbf{Bayesian inference, hierarchical models, and model misspecification.}
The fundamental difficulty in metabolic flux analysis (MFA) is that the
stoichiometric model $\dot{\mathbf{x}} = S\mathbf{v}$ is always misspecified:
unmeasured fluxes, unknown regulatory interactions, compartment-averaging
artefacts, and isotope label dilution all introduce systematic deviations
between model and data that point estimates cannot detect.
A Bayesian hierarchical approach handles misspecification by placing priors
over model parameters \emph{and} over the degree of model inadequacy: the
posterior distribution over flux vectors $p(\mathbf{v} \mid \mathbf{m})$
captures both estimation uncertainty and structural uncertainty jointly.

I have built and applied hierarchical Bayesian models in two directly relevant
contexts. The \href{https://github.com/fullscreen-triangle/gospel}{Gospel}
genomic analysis framework implements Bayesian network inference for
multi-omics data integration, using variational Bayes for posterior
approximation in models with $\mathcal{O}(10^6)$ latent variables.
The \href{https://github.com/fullscreen-triangle/syndrome}{Syndrome} disease
modelling framework uses Bayesian P(R\,|\,Q) --- the posterior probability of
answer quality given query class --- as a gating mechanism over an
eight-dimensional latent state, with continuous shrinkage priors over the
coherence indices $\eta_i \in [0,1]$.
For the MFA context, these translate directly: the latent flux vector
$\mathbf{v}$ plays the role of the disease state $\mathbf{D}$; the
isotopomer balance equations play the role of the stoichiometric network;
and the hierarchical prior over flux magnitudes encodes the biological
constraints (thermodynamic feasibility, capacity constraints, mass balance)
analogously to the sparse $\ell_1$ regularisation I use for minimum
intervention selection.

Beyond methodology, I have a theoretical result directly applicable to the
model misspecification problem. The \textbf{Backward Training Principle}
(developed in the \href{https://github.com/fullscreen-triangle/purpose}{Purpose}
framework) proves that a model calibrated at the worst-case misspecification
scenario --- the hardest, most constrained instance of the inference problem
--- is calibrated at all less-misspecified instances by feasibility inclusion:
\[
\mathcal{C}_{\mathrm{hard}} \subseteq \mathcal{C}_{\mathrm{easy}}
\implies
\epsilon\text{-competence on } \mathcal{C}_{\mathrm{hard}}
\implies
\epsilon\text{-competence on } \mathcal{C}_{\mathrm{easy}}.
\]
For a Bayesian MFA framework, the worst-case misspecification is the fully
unidentified model (maximum parameter uncertainty, minimum data information);
calibration there propagates to all better-identified sub-models.
A Forward Training Collapse Theorem establishes a positive Vapnik-Chervonenkis
failure gap for the conventional easy-to-hard approach; a Sample Complexity
Theorem gives $(N{+}1)\times$ savings for a cascade of depth $N$.
This is a formal contribution to the project's goal of designing hierarchical
models that capture misspecification robustly.

\textbf{Differentiable programming and HPC-optimised sampling.}
The project calls for differentiable, scalable inference algorithms using
automatic differentiation and HPC-optimised sampling in Python and C++.
My implementations use PyTorch as the primary differentiable computing
substrate: every model in my frameworks is built on \texttt{torch.autograd}
so that gradients through the full model --- including through stochastic
nodes via the reparameterisation trick --- are available for gradient-based
MCMC (NUTS, HMC) and variational inference (ELBO optimisation).
For HPC scale, I implement distributed sampling via
\href{https://github.com/ray-project/ray}{Ray}: multiple Markov chains run
in parallel across cluster nodes, with shared memory for the likelihood
evaluation and asynchronous parameter collection.
For performance-critical likelihood evaluations, I implement in Rust and
expose to Python via PyO3 bindings, achieving 15--25$\times$ speedup over
pure Python for the inner loops of MCMC sampling.
The \href{https://github.com/fullscreen-triangle/four-sided-triangle}{Four-Sided-Triangle}
framework demonstrates this architecture in a deployed system: 8-stage
inference pipeline with Rust core, Ray-distributed parallel execution,
and Python orchestration layer.

\textbf{Metabolic measurement data and flux analysis background.}
My doctoral research at TUM / Bayerisches Forschungsinstitut (2019--2021)
was in Computational Lipidomics: mass spectrometry-based quantification of
lipid species across hundreds of samples from distributed clinical sites.
Lipidomics is the downstream measurement of lipid metabolic fluxes --- the
lipid species abundances in a sample are the integrated outputs of the
synthesis, modification, transport, and degradation reactions in the lipid
metabolic network.
The pipeline I worked on covered: peak detection and alignment across
mzML files $>$100\,GB, ion annotation against LIPID MAPS and LipidBlast,
multi-omics integration combining lipidomics with clinical covariates, and
federated statistical modelling across distributed sites.
This gives me direct familiarity with the data structures and measurement
uncertainties (isotopologue distributions, matrix effects, ion suppression)
that a probabilistic MFA framework must handle.
The \href{https://github.com/fullscreen-triangle/lavoisier}{Lavoisier}
framework I subsequently built automates this pipeline at production scale:
dual numerical and CNN-based spectral analysis, 98.9\,\% NIST accuracy,
Ray-distributed for HPC, $>$100\,GB mzML support.
Isotopically labelled metabolomics data (the primary input to $^{13}$C-MFA)
is structurally identical to the lipidomics data I have processed: the
difference is the labelling scheme applied to the sample, not the measurement
instrument or data format.

\textbf{Open-source contributions and collaborative research.}
All frameworks I have built are publicly available on GitHub with documentation,
validation reports, and reproducibility instructions. I contribute to
open-source scientific computing projects and have implemented library-level
components (Rust statistical kernels, Python API wrappers for genomic
databases, WebGPU compute shaders for physical simulation) that are used
in the deployed applications.

\textbf{Background.}
I hold an MSc in Computational Biology (Universität Tübingen) and conducted
doctoral research in Computational Lipidomics at TUM. I work primarily in
Python and Rust; C experience from Fraunhofer IGB and scientific computing
contexts. I am legally resident in Germany with full work authorisation and
available immediately; willing to relocate to Jülich.
English is native; German is conversational.

\vspace{1.5em}
Yours sincerely,

\vspace{2.5em}
\textbf{Kundai Farai Sachikonye}

\end{document}
